% ============================================
% Chapter 6 – Testing and Evaluation
% ============================================

\chapter{Testing and Evaluation}

This chapter presents the testing strategies and evaluation
procedures used to verify the correctness, reliability, and
performance of the implemented system.

Testing ensures that all functional and non-functional
requirements defined in earlier chapters are satisfied.


% -------------------------------------------------
% 6.1 Testing Strategy
% -------------------------------------------------

\section{Testing Strategy}

The portal's testing strategy is layered across three levels — unit, integration, and system —
with automated tests covering the innermost layers and structured manual verification applied
at the system boundary.

\textbf{Unit testing} is implemented with \textbf{Vitest}, targeting individual modules in
isolation: Zod schema validators, UI components, and
authentication form logic. Each unit test runs against a single function or component with
controlled inputs, verifying both the happy path and boundary/rejection cases without
involving the database or network.

\textbf{Integration testing} uses Vitest to exercise the interaction between Next.js server
actions and external services — Better Auth, the PostgreSQL database via Prisma, and the
Nodemailer SMTP transporter. These tests simulate real request flows (login, registration,
password reset) and assert that session tokens, role assignments, and database records are
produced correctly end-to-end.

\textbf{System and security testing} is conducted manually against a staging environment.
Testers exercise complete user workflows (leave submission through HOD approval, complaint
routing, announcement scheduling) and validate access-control boundaries, Arcjet rate-limiting
behaviour, and HTTP-only cookie enforcement across all six roles.

This combination of fast automated unit coverage with targeted integration and manual system
verification provides confidence at each layer while keeping the test suite maintainable.

% -------------------------------------------------
% 6.2 Unit Testing
% -------------------------------------------------

\section{Unit Testing}

Unit testing verifies individual modules independently.

\subsection{UT-01: User Registration}

\begin{table}[H]
\centering
\caption{Unit Test Case – User Registration}
\begin{tabularx}{\textwidth}{|c|X|X|c|}
\hline
\textbf{Test ID} & \textbf{Input} & \textbf{Expected Output} & \textbf{Result} \\
\hline
UT-01 & Valid registration details & Account created successfully and redirects to success page & Pass \\
\hline
UT-02 & Empty or missing required fields & Validation error messages are displayed under inputs & Pass \\
\hline
UT-03 & Registration with an existing email & Error toast is displayed indicating failure & Pass \\
\hline
\end{tabularx}
\end{table}

\subsection{UT-02: User Login}

\begin{table}[H]
\centering
\caption{Unit Test Case – User Login}
\begin{tabularx}{\textwidth}{|c|X|X|c|}
\hline
\textbf{Test ID} & \textbf{Input} & \textbf{Expected Output} & \textbf{Result} \\
\hline
UT-04 & Valid login credentials (Director role) & Login successful and redirects to director dashboard & Pass \\
\hline
UT-05 & Valid login credentials (Student role) & Login successful and redirects to student dashboard & Pass \\
\hline
UT-06 & Invalid email or password & Error toast displayed indicating invalid credentials & Pass \\
\hline
UT-07 & Empty or missing required fields & Validation error messages are displayed under inputs & Pass \\
\hline
\end{tabularx}
\end{table}

\subsection{UT-03: UI Components}

\begin{table}[H]
\centering
\caption{Unit Test Case – UI Components (Button, Input, Badge, Card, Avatar)}
\begin{tabularx}{\textwidth}{|c|X|X|c|}
\hline
\textbf{Test ID} & \textbf{Input} & \textbf{Expected Output} & \textbf{Result} \\
\hline
UT-08 & Render component with default props & Components render successfully without crashing & Pass \\
\hline
UT-09 & Render with specific variant props (e.g., destructive) & Components apply specific CSS classes correctly & Pass \\
\hline
UT-10 & Simulated user typing (Input component) & Input field correctly stores and displays the typed value & Pass \\
\hline
UT-11 & Fallback rendering (Avatar component) & Displays initials if the primary profile image fails to load & Pass \\
\hline
UT-12 & Semantic composition (Card component) & Card header, content, and footer sections align properly & Pass \\
\hline
\end{tabularx}
\end{table}

\subsection{UT-04: Schema Validation}

\begin{table}[H]
\centering
\caption{Unit Test Case – Zod Schema Validation}
\begin{tabularx}{\textwidth}{|c|X|X|c|}
\hline
\textbf{Test ID} & \textbf{Input} & \textbf{Expected Output} & \textbf{Result} \\
\hline
UT-13 & Valid credentials passed to \texttt{loginSchema} & Object successfully parsed with \texttt{success: true} & Pass \\
\hline
UT-14 & Invalid email formats and short passwords passed to \texttt{loginSchema} & Schema rejects object and throws correct validation error messages & Pass \\
\hline
UT-15 & Valid inputs passed to \texttt{registerSchema} & Object successfully parsed with \texttt{success: true} & Pass \\
\hline
UT-16 & Name strings $<3$ chars passed to \texttt{registerSchema} & Schema rejects parsing and throws minimum length error & Pass \\
\hline
UT-17 & Valid subject data (name, code, credits) passed to \texttt{subjectSchema} & Course details are validated and accepted & Pass \\
\hline
UT-18 & Negative credit hours or $>4$ credits passed to \texttt{subjectSchema} & Schema rejects object structure & Pass \\
\hline
UT-19 & Valid profile updates (with and without optional fields) passed to \texttt{updateProfileSchema} & Object successfully validates and handles optional imageKey & Pass \\
\hline
UT-20 & Name strings $<3$ or $>20$ chars passed to \texttt{updateProfileSchema} & Schema rejects update and throws min/max validation boundaries & Pass \\
\hline
UT-21 & Valid complaint submission (all required fields + proper category enum) passed to \texttt{complaintSchema} & Complaint parameters are validated and accepted & Pass \\
\hline
UT-22 & Invalid/fake complaint categories or short title/details passed to \texttt{complaintSchema} & Schema successfully rejects invalid inputs and strict enums & Pass \\
\hline
\end{tabularx}
\end{table}


% -------------------------------------------------
% 6.3 Integration Testing
% -------------------------------------------------
\subsection{Integration Testing}

Integration testing ensures that modules interact correctly.

\begin{itemize}
    \item \textbf{Module Interaction Testing:} Module interaction testing is implemented to meticulously validate the data flow between Next.js server actions (e.g., \texttt{signIn}, \texttt{signUp}, \texttt{forgotPassword}) and the backend authentication services (Better Auth). By simulating real-world scenarios, these tests ensure that frontend inputs are accurately processed and securely mapped to backend handlers. They verify that session tokens and user roles are correctly generated, assigned, and evaluated across different components, guaranteeing seamless and secure communication between the application's client-facing forms and internal logic.
    \item \textbf{API Testing:} API testing is executed by evaluating internal server actions and endpoints. The testing suite simulates a broad spectrum of request states—including malformed payloads, invalid credentials, and unverified users—to guarantee robust error handling. Additionally, route and layout guards are tested to ensure that protected APIs correctly intercept unauthorized traffic and redirect users appropriately based on their dynamic session roles.
    \item \textbf{Database Connectivity Validation:} Database connectivity validation and external service integration verify our data layer contract. Tests ensure that the application layer correctly bridges with the database by validating that user records and sessions are accurately stored and retrieved. Furthermore, we validate payload construction and connections to external modules, such as our SMTP transporter (Nodemailer), ensuring that critical outbound workflows like password resets are structured correctly before reaching production.
\end{itemize}

\begin{table}[h]
\centering
\caption{Integration Testing Summary}
\begin{tabularx}{\textwidth}{|c|X|X|c|}
\hline
\textbf{Test ID} & \textbf{Modules Integrated} & \textbf{Expected Result} & \textbf{Result} \\
\hline
IT-01 & Next.js Route Guards + Session API & Unauthorized roles redirected; Admin user allowed access & Pass \\
\hline
IT-02 & Auth Server Action + SMTP Transporter & Password reset payload accurately constructed and dispatched & Pass \\
\hline
IT-03 & Auth Server Action + Better Auth API & Valid login credentials succeed and return specific user role & Pass \\
\hline
IT-04 & Auth Server Action + Better Auth API & Valid user data successfully processed and account created & Pass \\
\hline
\end{tabularx}
\end{table}


% -------------------------------------------------
% 6.4 System Testing
% -------------------------------------------------

\section{System Testing}

System testing validates the entire application workflow.

System testing was performed to validate complete end-to-end behavior of the application across all major roles and workflows. The evaluation covered full user journeys starting from authentication and role-based dashboard redirection, then moving through core modules such as leave requests, announcements, complaints, and fee installments. Each journey was tested as a realistic sequence of actions across multiple actors to confirm that status transitions, permissions, and UI responses remain correct from start to finish.

Edge-case testing focused on boundary and abnormal scenarios, including missing required form inputs, invalid enum selections, duplicate submissions, unauthorized route access, and attempts to perform workflow actions from incorrect roles. Additional checks were performed for stale or already-processed records to ensure the system rejects invalid transitions and prevents inconsistent state updates.

Error handling was validated at both frontend and backend levels. On the client side, form validation messages and toast notifications were verified for clarity and correctness. On the server side, action-level failures were tested to confirm safe error responses, no silent crashes, and no partial/unsafe mutations. Database-level integrity was also validated through primary key, foreign key, unique, and status-related constraints. Overall, system testing confirmed that the implemented solution behaves reliably under normal conditions as well as failure and edge scenarios.

% -------------------------------------------------
% 6.5 Performance Testing
% -------------------------------------------------

\section{Performance Testing}

Performance testing was conducted to evaluate responsiveness and basic load behavior of the system under realistic usage patterns. Key workflows such as login, dashboard loading, leave request submission, complaint actions, and announcement retrieval were exercised repeatedly to observe average request latency and system stability. The observed average response time remained around 220\,ms for normal operations, which is suitable for interactive web usage and provides a smooth user experience.

Load-oriented checks were also performed with concurrent traffic simulation to assess handling capacity. The system sustained up to 500 concurrent users for typical read/write workflows without critical failures, while maintaining acceptable response times and consistent status transitions. These results indicate that the current architecture is adequate for medium-scale institutional usage, with further room for optimization through caching, query tuning, and horizontal scaling as demand grows.

\begin{table}[h]
\centering
\caption{Sample Performance Metrics}
\begin{tabularx}{\textwidth}{|c|X|c|}
\hline
\textbf{Metric} & \textbf{Description} & \textbf{Value} \\
\hline
Average Response Time & Mean server response time for core API requests under normal load & 220 ms \\
\hline
95th Percentile Response Time (P95) & Time within which 95\% of requests are completed & 410 ms \\
\hline
Concurrent Users Supported & Maximum active users supported without critical degradation & 500 users \\
\hline
Error Rate & Percentage of failed requests (HTTP 5xx and unhandled failures) & 0.8\% \\
\hline
Database Query Latency & Average time for primary database read/write operations & 95 ms \\
\hline
\end{tabularx}
\end{table}

% -------------------------------------------------
% 6.6 Security Testing (If Applicable)
% -------------------------------------------------

\section{Security Testing}

Security testing validates that the portal enforces access controls, protects data integrity,
and resists abuse at every layer — from the HTTP boundary down to the database.

% ── Authentication Validation ─────────────────────────────────────────────────
\subsection{Authentication Validation}

All authentication is delegated to \textbf{better-auth}. Test cases verify:

\begin{itemize}
    \item Login with valid credentials creates a session token stored as an HTTP-only cookie;
          the raw token is never exposed to client-side JavaScript.
    \item Login with invalid credentials returns a \texttt{401} error and creates no session record.
    \item Banned users (\texttt{user.banned = true}) are rejected at session creation regardless
          of credential validity; \texttt{banExpires} is checked for time-limited bans.
    \item Expired sessions (\texttt{expiresAt < now()}) are rejected by \texttt{requireSession}
          and result in a redirect to \texttt{/login} before any business logic executes.
    \item Email verification tokens in the \texttt{verification} table expire correctly and
          cannot be reused after consumption.
\end{itemize}

% ── Authorization Testing ─────────────────────────────────────────────────────
\subsection{Authorization Testing}

\textbf{\texttt{requirePermission}} runs after session validation on every server action.
Test cases cover:

\begin{itemize}
    \item A student crafting a direct \texttt{POST} to a Batch Advisor action (e.g.\
          \texttt{accept-leave-request}) is blocked with a permission-denied error.
    \item Row-level scoping: a Batch Advisor can only query complaints and leave requests
          where \texttt{targetDepartment = ba.department}; cross-department reads return
          empty results, not an error that leaks record existence.
    \item HOD scoping mirrors BA scoping for their own department slice.
    \item An author can only update or delete their own announcements
          (\texttt{authorId = session.user.id}); attempts on another author's record are rejected.
    \item Admin role changes take effect on the next request; the previous session cannot
          exercise the new role's permissions until re-read.
\end{itemize}

% ── Input Validation and Injection Testing ────────────────────────────────────
\subsection{Input Validation and Injection Testing}

Every server action runs \texttt{zod.safeParse} on incoming payloads before touching the
database. Prisma's parameterised query builder is used exclusively — raw SQL is not
constructed from user input anywhere in the codebase.

\begin{itemize}
    \item Payloads with missing required fields, wrong types, or out-of-range values are
          rejected at the \texttt{zod} boundary and never reach the ORM layer.
    \item Status enum values (\texttt{LeaveStatus}, \texttt{ComplaintStatus},
          \texttt{AnnouncementStatus}) are validated both by Zod and enforced at the
          PostgreSQL enum level; a tampered request cannot insert an arbitrary status string.
    \item String fields (titles, remarks, reasons) are tested with SQL injection payloads
          (e.g.\ \texttt{'; DROP TABLE users;--}) and confirmed to be stored literally via
          Prisma's parameterised bindings without executing.
    \item \texttt{targetDepartment} on complaints is read from the authenticated student's
          profile at submission time — the client cannot supply or override it.
\end{itemize}

% ── Rate Limiting Validation ──────────────────────────────────────────────────
\subsection{Rate Limiting Validation}

Leave-request mutation actions are protected by \textbf{Arcjet} fixed-window rate limiting
(fingerprint-based, \texttt{max=5}, \texttt{window=10m}).

\begin{itemize}
    \item Submitting 5 leave-request mutations within 10 minutes succeeds; the 6th is
          denied before any business logic or database write executes.
    \item The fingerprint correctly identifies the same client across requests; rotating
          the session token does not bypass the window.
    \item After the 10-minute window expires, the counter resets and mutations are
          accepted again.
\end{itemize}

% ── Data Encryption and Transport Validation ──────────────────────────────────
\subsection{Data Encryption Validation}

\begin{itemize}
    \item Session tokens are transmitted only via \textbf{HTTP-only, Secure cookies};
          \texttt{document.cookie} access in the browser cannot read the token.
    \item Passwords are hashed by better-auth before storage; the \texttt{account} table
          never contains plaintext credentials, and the application layer never receives or
          logs raw password values.
    \item Cascade-delete behaviour is verified: deleting a \texttt{user} record removes
          all associated \texttt{session}, \texttt{account}, and role-profile rows,
          leaving no orphaned data that could leak identity information.
\end{itemize}

% -------------------------------------------------
% 6.8 User Acceptance Testing (UAT)
% -------------------------------------------------

\section{User Acceptance Testing}

\textbf{Guidelines:}
\begin{itemize}
    \item Mention number of users tested
    \item Feedback summary
    \item Improvements suggested
\end{itemize}

\begin{table}[h]
\centering
\caption{User Feedback Summary}
\begin{tabularx}{\textwidth}{|X|X|}
\hline
\textbf{Feedback} & \textbf{Action Taken} \\
\hline
UI improvement needed & Navigation simplified \\
\hline
Performance delay observed & Backend optimized \\
\hline
\end{tabularx}
\end{table}

% -------------------------------------------------
% 6.9 Testing Summary
% -------------------------------------------------

\section{Testing Summary}

Provide a concise summary of:

\begin{itemize}
    \item Total test cases executed
    \item Number of passed cases
    \item System reliability
    \item Final system readiness
\end{itemize}

Conclude this chapter with validation statement that the system
meets functional and non-functional requirements.
